% Teil II: der durchgerechnete Weg.

\subsection*{Gegeben und gesucht}
Gegeben: Klotz $m$, Breite $B=\qty{0.90}{m}$, Gleitreibungszahl
$\mu=\qty{0.34}{}$, Schubkraft $F$ auf Höhe $h$, konstante
Gleitgeschwindigkeit, $g=\qty{9.81}{m/s^2}$.
Gesucht: maximale Höhe $h_\mathrm{max}$, bei der der Klotz nicht kippt.

\subsection*{Ansatz}
Aus dem Kräftegleichgewicht (konstante Geschwindigkeit) $F=\mu mg$. Drehmomente
um den Kipp-Punkt $P$ an der vorderen unteren Kante; Bedingung "nicht
kippen" $Fh\le mg\,\frac B2$.

\subsection*{Rechnung}
Horizontalkräfte im Gleichgewicht:
\begin{equation}
  F = \mu\, m g.
\end{equation}
Moment um den Kipp-Punkt $P$ (vordere untere Kante). Die Reibungskraft wirkt
kurz vor dem Kippen \emph{an $P$} selbst und erzeugt dadurch hier kein Moment;
das Gewicht wirkt im Schwerpunkt am Hebelarm $B/2$ (stabilisierend). Maximal
darf das kippende Moment des Schubs gerade gleich sein:
\begin{equation}
  F\cdot h = m g\cdot \frac B2.
\end{equation}
Einsetzen von $F=\mu m g$ ($m g$ kürzt sich):
\begin{equation}
  \mu\,m g\, h = m g\,\frac B2
  \;\Rightarrow\;
  h = \frac{B}{2\mu}
  = \frac{0.90}{2\cdot 0.34}\ \si{m}
  = \qty{1.32}{m}.
\end{equation}

\textbf{Antwort zur Reibungskraft:} Konstruktiv wirkt die Gleitreibung an der
gesamten Kontaktfläche; kurz \emph{vor} dem Kippen (wenn der Klotz nur noch um
$P$ drehen möchte) wandert die Normalkraft samt Reibung zum Kipp-Punkt $P$.
Dadurch beträgt der Hebelarm der Reibungskraft um $P$ null und ihr Moment
verschwindet.

\begin{ergebnis}
\begin{equation*}
  h_\mathrm{max} = \frac{B}{2\mu} = \qty{1.32}{m}.
\end{equation*}
Für $h\le\qty{1.32}{m}$ (und $h\le H$) kippt der Klotz nicht; darüber kippt er
trotz Gleitens um die vordere Kante.
\end{ergebnis}

\begin{kontrolle}
\textbf{Einheiten:} $B/(2\mu)$: $m/\mathrm{dimensionslos}=m$. \checkmark

\textbf{Grenzfälle:} Kleines $\mu$ (glatter Boden) ermöglicht eine grosse
Höhe $h_\mathrm{max}$; für $\mu\to0$ kippt der Klotz nie (er rutscht nur).
Grosses $B$ (breiter Klotz) vergrössert ebenfalls $h_\mathrm{max}$.
Plausibel.

\textbf{Konsistenz:} eine höhere Angriffshöhe $h$ vergrössert
das kippende Moment $Fh$; das Reibungsgleichgewicht $F=\mu mg$ ist dabei
notwendig (bei gegebener Reibung nicht willkürlich).
\end{kontrolle}

\begin{merksatz}
Kippen untersucht man mit Drehmomenten um die Kippkante $P$: stabilisierend
wirkt das Gewicht ($mg$ bei Hebelarm $B/2$), kippend die Schubkraft ($F$ bei
Hebelarm $h$). Zusammen mit dem Reibungsgleichgewicht $F=\mu mg$ ergibt sich
$h_\mathrm{max}=B/(2\mu)$. Die Kipppunkt-Reibkraft trägt zum Moment nichts bei.
\end{merksatz}